\section{Miaoya: un producto de internet que usa bien la IA}

Esta sección pasa de la trayectoria personal a un caso de producto. Entender Miaoya (妙鸭) es la primera puerta para entender toda la entrevista: es el caso de éxito más conocido de Zhang Yueguang (张月光) y también el contraejemplo que más tarde usó para reflexionar sobre «qué es realmente AI Native».

El éxito de Miaoya proviene de varias decisiones de producto: la tecnología de imagen ya era relativamente madura en ese momento; no hacer generación de imágenes de propósito general al estilo de Midjourney; elegir un nicho suficientemente vertical, el retrato fotográfico; evitar la larga cadena del comercio electrónico y optar por un negocio de retratos dirigido al consumidor (2C) al que se le puede poner precio; y reducir el objetivo de resultado a «realista, parecido, bello». Todas estas decisiones son excelentes, pero Zhang Yueguang concluyó después que Miaoya seguía respondiendo a una mentalidad de producto de internet.

\begin{figure}[H]
\centering
\includegraphics[width=0.94\textwidth]{miaoya-not-ai-native.png}
\caption{Por qué Miaoya no es AI Native: entrada fija, cadena de procesamiento fija, salida fija.}
\end{figure}

\begin{knowledgebox}{Lectura de la figura: la estabilidad de Miaoya proviene de limitar los grados de libertad}
Los tres segmentos de la figura son entrada fija, cadena de procesamiento fija y salida fija. El usuario sube unas 20 fotos, el sistema entrena un modelo personal y luego produce los retratos mediante plantillas fijas y una cadena de generación fija. Usa muy bien la capacidad de la IA, pero sigue obteniendo resultados estables a cambio de limitar la libertad del usuario.
\end{knowledgebox}

\subsection{Realista, parecido, bello}

Antes se dijo que la cadena de procesamiento de Miaoya es fija, pero una cadena fija no equivale a una cadena sencilla. Su verdadera dificultad está en lograr, para un escenario muy acotado, una función de efecto lo bastante sólida.

El estándar de resultado de Miaoya consta de tres elementos. Realista: que no se note que lo generó una IA; parecido: que se parezca lo suficiente al propio usuario, de modo que lo que se comparte tenga que ver con «yo»; bello: que se vea un poco mejor que uno mismo en la realidad. Este triángulo explica por qué se volvió un éxito masivo: no exhibe la tecnología de IA, sino que la convierte en una experiencia de consumo intensa.

\begin{importantbox}{La lección de producto de Miaoya}
Un éxito masivo de IA no necesariamente proviene del modelo más general; puede provenir de una función de efecto bien definida dentro de un escenario vertical. La función de efecto de Miaoya es justamente «realista, parecido, bello».
\end{importantbox}

\subsection{Por qué sigue siendo un producto de internet}

La función de efecto explica por qué Miaoya tuvo éxito; esta sección explica por qué, aun así, no es AI Native. Lo decisivo no es si el resultado deslumbra, sino si el producto sigue dependiendo de un flujo predefinido.

En esencia, Miaoya sigue siendo un diseño de flujo: el usuario sube fotos siguiendo pasos predefinidos, elige una plantilla y el sistema genera el resultado a lo largo de una cadena de proceso fija. La entrada del usuario no es abierta y la salida tampoco. El gerente de producto puede diseñar el flujo completo, los ingenieros lo implementan tal como está definido y la experiencia es controlable. Es un triunfo de la metodología de producto de internet, no el despliegue completo de la metodología AI Native.

\subsection{Resumen de la sección}

Miaoya muestra que usar bien la capacidad de la IA no equivale a que la IA haya reescrito el paradigma de producto. Es un excelente «producto de internet con IA añadida», pero no es AI Native en el sentido en que Zhang Yueguang lo entendió después.
